\documentclass[10pt,a4paper]{amsart}
\usepackage{amsmath,amssymb,amsthm,graphicx}
% Retained minimal preamble for the Stacks Project chapter more-algebra.tex (More on Algebra).
% The chapter file pinned-original-source.tex carries only \input{preamble} on its first line, so no
% preamble text lives inside the chapter itself.  The Stacks Project uses one single shared preamble for
% every chapter; the snapshot stages it identically as sources/stacks-chow/preamble.tex and
% sources/stacks-etale-cohomology/preamble.tex.  The block below is copied mechanically, line for line,
% from sources/stacks-chow/preamble.tex lines 14-18 (the Xy-pic setup) and 154-217 (the native theorem
% declarations together with the mathematical macros).  Nothing was added, removed or redefined; the
% wrapper class and the amsmath/amssymb/amsthm/graphicx base packages are supplied by the assembler.
% Chapter: more-algebra.tex -- More on Algebra.
% For commutative diagrams we use Xy-pic
\usepackage[all]{xy}

% We use 2cell for 2-commutative diagrams.
\xyoption{2cell}

% Theorem environments.
%
\theoremstyle{plain}
\newtheorem{theorem}[subsection]{Theorem}
\newtheorem{proposition}[subsection]{Proposition}
\newtheorem{lemma}[subsection]{Lemma}

\theoremstyle{definition}
\newtheorem{definition}[subsection]{Definition}
\newtheorem{example}[subsection]{Example}
\newtheorem{exercise}[subsection]{Exercise}
\newtheorem{situation}[subsection]{Situation}

\theoremstyle{remark}
\newtheorem{remark}[subsection]{Remark}
\newtheorem{remarks}[subsection]{Remarks}

\numberwithin{equation}{subsection}

% Macros
%
\def\lim{\mathop{\mathrm{lim}}\nolimits}
\def\colim{\mathop{\mathrm{colim}}\nolimits}
\def\Spec{\mathop{\mathrm{Spec}}}
\def\Hom{\mathop{\mathrm{Hom}}\nolimits}
\def\Ext{\mathop{\mathrm{Ext}}\nolimits}
\def\SheafHom{\mathop{\mathcal{H}\!\mathit{om}}\nolimits}
\def\SheafExt{\mathop{\mathcal{E}\!\mathit{xt}}\nolimits}
\def\Sch{\mathit{Sch}}
\def\Mor{\mathop{\mathrm{Mor}}\nolimits}
\def\Ob{\mathop{\mathrm{Ob}}\nolimits}
\def\Sh{\mathop{\mathit{Sh}}\nolimits}
\def\NL{\mathop{N\!L}\nolimits}
\def\CH{\mathop{\mathrm{CH}}\nolimits}
\def\proetale{{pro\text{-}\acute{e}tale}}
\def\etale{{\acute{e}tale}}
\def\QCoh{\mathit{QCoh}}
\def\Ker{\mathop{\mathrm{Ker}}}
\def\Im{\mathop{\mathrm{Im}}}
\def\Coker{\mathop{\mathrm{Coker}}}
\def\Coim{\mathop{\mathrm{Coim}}}

% Boxtimes
%
\DeclareMathSymbol{\boxtimes}{\mathbin}{AMSa}{"02}

%
% Macros for moduli stacks/spaces
%
\def\QCohstack{\mathcal{QC}\!\mathit{oh}}
\def\Cohstack{\mathcal{C}\!\mathit{oh}}
\def\Spacesstack{\mathcal{S}\!\mathit{paces}}
\def\Quotfunctor{\mathrm{Quot}}
\def\Hilbfunctor{\mathrm{Hilb}}
\def\Curvesstack{\mathcal{C}\!\mathit{urves}}
\def\Polarizedstack{\mathcal{P}\!\mathit{olarized}}
\def\Complexesstack{\mathcal{C}\!\mathit{omplexes}}
% \Pic is the operator that assigns to X its picard group, usage \Pic(X)
% \Picardstack_{X/B} denotes the Picard stack of X over B
% \Picardfunctor_{X/B} denotes the Picard functor of X over B
\def\Pic{\mathop{\mathrm{Pic}}\nolimits}
\def\Picardstack{\mathcal{P}\!\mathit{ic}}
\def\Picardfunctor{\mathrm{Pic}}
\def\Deformationcategory{\mathcal{D}\!\mathit{ef}}

% --- appended 2026-10-12 (Chief): the project own definitions of the reference/slogan/history
% display environments, copied mechanically from the shared Stacks preamble lines 8-12
% (\usepackage{verbatim} plus the three \newenvironment lines, which the chapter preamble keeps
% verbatim from the project preamble). Nothing added, removed or redefined: these environments
% carry no mathematical content and the project itself typesets them as comments. Without these
% three lines a counted statement whose own text uses them fails with Environment slogan undefined.
\usepackage{verbatim}
\newenvironment{reference}{\comment}{\endcomment}
\newenvironment{slogan}{\comment}{\endcomment}
\newenvironment{history}{\comment}{\endcomment}
\usepackage{cleveref}
\crefname{originallemma}{original lemma}{original lemmas}
\Crefname{originallemma}{Original lemma}{Original lemmas}
\makeatletter
\newlabel{lemma-quotient-by-idempotent}{{0ALR}{source}{}{sourcepointer-1}{}}
\newlabel{lemma-quotient-by-idempotent@cref}{{[originallemma][0][]0ALR}{[1][0][]source}{}{}{}}
\makeatother
\crefname{originallemma}{original lemma}{original lemmas}
\Crefname{originallemma}{Original lemma}{Original lemmas}
\makeatletter
\newlabel{algebra-lemma-dimension-formula}{{02IJ}{source}{}{sourcepointer-2}{}}
\newlabel{algebra-lemma-dimension-formula@cref}{{[originallemma][0][]02IJ}{[1][0][]source}{}{}{}}
\makeatother
\crefname{originallemma}{original lemma}{original lemmas}
\Crefname{originallemma}{Original lemma}{Original lemmas}
\makeatletter
\newlabel{algebra-lemma-finite-after-completion}{{0394}{source}{}{sourcepointer-3}{}}
\newlabel{algebra-lemma-finite-after-completion@cref}{{[originallemma][0][]0394}{[1][0][]source}{}{}{}}
\makeatother
\crefname{originallemma}{original lemma}{original lemmas}
\Crefname{originallemma}{Original lemma}{Original lemmas}
\makeatletter
\newlabel{algebra-lemma-silly}{{00DS}{source}{}{sourcepointer-4}{}}
\newlabel{algebra-lemma-silly@cref}{{[originallemma][0][]00DS}{[1][0][]source}{}{}{}}
\makeatother
\title{A Noetherian local ring that is not formally catenary is not universally catenary}
\author{The Stacks Project Authors}
\date{stacks-project:more-algebra.tex:local-snapshot-20261012; GNU FDL 1.2 or later}
\begin{document}
\maketitle
\setcounter{subsection}{1}
\noindent\textit{Editorial note.} All delivered mathematical text is the author's own, copied from the pinned Stacks Project chapter file more-algebra.tex; no statement or proof step was rewritten, reordered or removed. The counted span is the author environment and its own proof environment. References inside the span are handed to the assembler as navigation stubs that carry the canonical tag of the cited tag only; whether each of them is load-bearing is recorded on the stub itself.\par\medskip
\noindent\textit{Allowed standard prerequisites (editorial).}\begin{itemize}
\item A finite type extension of the ring can be built whose completion is the quotient of the completion by a prescribed ideal that has only the maximal ideal as associated prime. invoked at proof lines 32335-32340. Produces the finite type ring B with completion A-hat modulo J in the case of dimension one.
\item The dimension formula for a flat local map of Noetherian local rings, comparing dimensions with the transcendence degrees of the residue field and of the fibre. invoked at proof lines 32348-32358. Supplies the inequality that contradicts the dimension count once a prime lying over the zero ideal is produced.
\item Prime avoidance lets one choose an element of the maximal ideal avoiding finitely many prescribed primes. invoked at proof lines 32378-32380. Produces the element f used in the induction step so that the quotient by f is again not formally catenary with the invariant decreased.
\end{itemize}\medskip
\ifdefined\hypertarget\hypertarget{sourcepointer-1}{}\fi
\noindent\textit{Source reference.} Original lemma 0ALR: sources/stacks-more-algebra/pinned-original-source.tex, source lines 32037--32037. Invoked at proof lines 32338-32340 to produce the finite type ring B whose completion is A-hat modulo a power of q; the verification that the maximal ideal is the only associated prime of that ideal, which is the hypothesis of the reference, is inside the counted span.\par\medskip
\ifdefined\hypertarget\hypertarget{sourcepointer-2}{}\fi
\noindent\textit{Source reference.} Original lemma 02IJ: sources/stacks-chow/tags.txt, source lines 1--1. Applied at proof lines 32350-32360 to derive an inequality contradicting the computed dimension of the localization of B once a prime over the zero ideal is available; producing that prime and the dimension computation stay inside the counted span.\par\medskip
\ifdefined\hypertarget\hypertarget{sourcepointer-3}{}\fi
\noindent\textit{Source reference.} Original lemma 0394: sources/stacks-chow/tags.txt, source lines 1--1. Used at proof lines 32341-32343 to identify the completion of B with the adic completion at the maximal ideal and to compute its dimension; the surrounding dimension count is inside the counted span.\par\medskip
\ifdefined\hypertarget\hypertarget{sourcepointer-4}{}\fi
\noindent\textit{Source reference.} Original lemma 00DS: sources/stacks-chow/tags.txt, source lines 1--1. Used at proof lines 32379-32380 to pick an element of the maximal ideal outside finitely many primes; the verification that the quotiented ring is again not formally catenary with the invariant decreased is inside the counted span.\par\medskip
\begin{lemma}
\label{lemma-not-formally-catenary}
Let $(A, \mathfrak m)$ be a Noetherian local ring which is not
formally catenary. Then $A$ is not universally catenary.
\end{lemma}

\begin{proof}
By assumption there exists a minimal prime $\mathfrak p \subset A$ such that
the spectrum of $A^\wedge /\mathfrak p A^\wedge$ is not equidimensional.
After replacing $A$ by $A/\mathfrak p$ we may assume that $A$
is a domain and that the spectrum of $A^\wedge$ is not equidimensional.
Let $\mathfrak q$ be a minimal prime of
$A^\wedge$ such that $d = \dim(A^\wedge/\mathfrak q)$
is minimal and hence $0 < d < \dim(A)$. We prove the lemma by induction
on $d$.

\medskip\noindent
The case $d = 1$. In this case $\dim(A^\wedge_\mathfrak q) = 0$.
Hence $A^\wedge_\mathfrak q$ is Artinian local and we see that
for some $n > 0$ the ideal $J = \mathfrak q^n$ maps to zero in
$A^\wedge_\mathfrak q$. It follows that $\mathfrak m$ is the
only associated prime of $J/J^2$, whence $\mathfrak m^m$ annihilates
$J/J^2$ for some $m > 0$. Thus we can use
Lemma \ref{lemma-quotient-by-idempotent}
to find $A \to B$ of finite type such that $B^\wedge \cong A^\wedge/J$.
It follows that $\mathfrak m_B = \sqrt{\mathfrak mB}$ is a maximal
ideal with the same residue field as $\mathfrak m$ and $B^\wedge$
is the $\mathfrak m_B$-adic completion
(Algebra, Lemma \ref{algebra-lemma-finite-after-completion}).
Then
$$
\dim(B_{\mathfrak m_B}) = \dim(B^\wedge) = 1 = d.
$$
Since we have the factorization $A \to B \to A^\wedge/J$ the inverse image
of $\mathfrak q/J$ is a prime $\mathfrak q' \subset \mathfrak m_B$ lying
over $(0)$ in $A$. Thus, if $A$ were universally catenary, the dimension
formula (Algebra, Lemma \ref{algebra-lemma-dimension-formula}) would give
\begin{align*}
\dim(B_{\mathfrak m_B})
& \geq
\dim((B/\mathfrak q')_{\mathfrak m_B}) \\
& =
\dim(A) + \text{trdeg}_A(B/\mathfrak q') -
\text{trdeg}_{\kappa(\mathfrak m)}(\kappa(\mathfrak m_B)) \\
& =
\dim(A) + \text{trdeg}_A(B/\mathfrak q')
\end{align*}
This contradiction finishes the argument in case $d = 1$.

\medskip\noindent
Assume $d > 1$. Let $Z \subset \Spec(A^\wedge)$ be the union of
the irreducible components distinct from $V(\mathfrak q)$.
Let $\mathfrak r_1, \ldots, \mathfrak r_m \subset A^\wedge$
be the prime ideals corresponding to irreducible components of
$V(\mathfrak q) \cap Z$ of dimension $> 0$.
Choose $f \in \mathfrak m$, $f \not \in A \cap \mathfrak r_j$
using prime avoidance (Algebra, Lemma \ref{algebra-lemma-silly}).
Then $\dim(A/fA) = \dim(A) - 1$ and there is some irreducible
component of $V(\mathfrak q, f)$ of dimension $d - 1$.
Thus $A/fA$ is not formally catenary and the invariant $d$ has
decreased. By induction $A/fA$ is not universally catenary, hence
$A$ is not universally catenary.
\end{proof}

\end{document}
